%----------------------------------------------------------------------------------------
% REKENEN MET MATRICES. Alle opgaven zijn eigen werk.
%
% Opbouw: Start, leerstof (theorie, werkwijze of uitleg, voorbeeld, korte oefening), en op een
% nieuwe pagina de Opdrachten (door elkaar, stijgende moeilijkheid), Koppel en Test jezelf.
% Labels van dit hoofdstuk beginnen met h1-.
%----------------------------------------------------------------------------------------

\mychapter{1}{Rekenen met matrices}

\lkr{Leerplan Basis: keuzedoel K3 (bewerkingen met matrices: optelling, scalaire vermenigvuldiging, matrixvermenigvuldiging, machtsverheffing en transpositie). Het leerplan vraagt geen bewijzen; alleen de eenvoudige eigenschappen krijgen er een, op het niveau van $2 \times 2$-matrices. Het accent ligt op het product: wanneer het bestaat, hoe je het uitrekent en wat het betekent in een context. Het rekentoestel komt pas in het onderdeel over machten.}

Een tabel met getallen kom je overal tegen: een puntenlijst, een uurrooster, de verkoop van een winkel per dag. Laat je de opschriften weg en houd je alleen de getallen over, dan heb je een matrix. Dat lijkt een kleine stap, maar zo kun je met hele tabellen tegelijk rekenen.

Optellen en vermenigvuldigen met een getal gaan zoals je verwacht: element per element. Het product van twee matrices is minder voor de hand liggend, maar het is de bewerking die matrices zo nuttig maakt. Je ziet het al gebeuren als je een factuur opmaakt: aantallen maal prijzen, en alles opgeteld. In hoofdstuk 2 gebruiken we het product om te voorspellen hoe een situatie evolueert.

\begin{oefeningen}[Start]
\oef{h1-start-1}

\oef{h1-verken}
\end{oefeningen}

\lkr{Laat oefening \ref{oef:h1-verken} maken vóór de uitleg. Vraag a is de optelling, b en c zijn het product van een rij met een kolom. Vraag d is de kern: leerlingen formuleren zelf dat je de getallen van een rij combineert met die van een kolom.}

%========================================================================================
\section{Matrices}

\subsection{Wat is een matrix?}

De filmzaal \textit{Polyglot} toont films in de originele taal. De tabel geeft hoeveel tickets er in een weekend verkocht werden, per dag en per taal van de film.

\figuren{\begin{tabular}{l|cccc}
 & Frans & Engels & Spaans & Italiaans \\ \hline
vrijdag & 48 & 112 & 35 & 20 \\
zaterdag & 96 & 180 & 64 & 41 \\
zondag & 70 & 150 & 52 & 38
\end{tabular}\qquad\raisebox{-4.2mm}{$\longrightarrow$}\qquad
$T = \begin{mat} 48 & 112 & 35 & 20 \\ 96 & 180 & 64 & 41 \\ 70 & 150 & 52 & 38 \end{mat}$}

Rechts staan alleen nog de getallen, in dezelfde rijen en kolommen. Zolang je weet wat de rijen en de kolommen betekenen, verlies je niets. Om een getal aan te wijzen, zeg je in welke rij en in welke kolom het staat: eerst de rij, dan de kolom.

\begin{theorie}
\begin{minipage}[t]{.56\linewidth}
\vspace{0pt}
Een \textbf{\textcolor{accent}{matrix}} is een rechthoekig rooster van getallen tussen vierkante haken. De getallen heten de \textbf{\textcolor{accent}{elementen}}.

Een matrix met $m$ rijen en $n$ kolommen heeft \textbf{\textcolor{accent}{dimensie}} $\formule{m \times n}$.

Het element in rij $i$ en kolom $j$ van de matrix $A$ noteer je $\formule{a_{ij}}$.
\end{minipage}\hfill
\begin{minipage}[t]{.4\linewidth}
\vspace{0pt}
\begin{tikzpicture}[font=\small]
  \matrix[getallen,nodes={minimum width=8mm}] (m) {a_{11} \& a_{12} \& a_{13} \& a_{14} \\ a_{21} \& a_{22} \& a_{23} \& a_{24} \\ a_{31} \& a_{32} \& a_{33} \& a_{34} \\};
  \markeer[accent!14]{m-2-1}{m-2-4}
  \markeer[accent!14]{m-1-3}{m-3-3}
  \markeer[fel!60]{m-2-3}{m-2-3}
  \node[left=4mm of m-2-1,font=\scriptsize,text=accent] {rij 2};
  \node[above=2.5mm of m-1-3,font=\scriptsize,text=accent] {kolom 3};
  \node[below=1mm of m.south,font=\scriptsize] {dimensie $3 \times 4$};
\end{tikzpicture}
\end{minipage}
\end{theorie}

Een matrix krijgt een hoofdletter als naam, en haar elementen dezelfde letter als kleine letter. Bij een $3 \times 4$-matrix zeg je: \textit{drie op vier}.

\voorbeeld{}{Wat is de dimensie van de matrix $T$ van de filmzaal? Wat is $t_{23}$, en wat betekent dat getal? Welk element geeft het aantal tickets voor een Engelse film op zondag?\label{vb:h1-element}}{%
\uitleg{De dimensie tel je af: eerst de rijen, dan de kolommen. Voor de betekenis van een element kijk je in de tabel welke dag bij de rij hoort en welke taal bij de kolom.}
\begin{stappen}
3 \times 4 & Er zijn drie rijen (de dagen) en vier kolommen (de talen). \\[1ex]
t_{23} = 64 & Rij 2 is zaterdag, kolom 3 is Spaans: op zaterdag werden 64 tickets voor een Spaanse film verkocht. \\[1ex]
t_{32} = 150 & Zondag is rij 3, Engels is kolom 2. Let op de volgorde: $t_{23}$ is een ander element.
\end{stappen}}

\nujij{Welk element van $T$ geeft het aantal tickets voor een Italiaanse film op vrijdag? Hoeveel is dat?}{Vrijdag is rij 1, Italiaans kolom 4: $t_{14} = 20$.}

\begin{oefeningen}
\oef{h1-element}
\end{oefeningen}

\subsection{Bijzondere matrices}

Sommige vormen komen zo vaak voor dat ze een eigen naam hebben. In een vierkante matrix noemen we de elementen $a_{11}, a_{22}, a_{33}, \ldots$ samen de \textbf{hoofddiagonaal}: die loopt van linksboven naar rechtsonder.

\begin{theorie}
{\centering\arrayrulecolor{accent}\renewcommand{\arraystretch}{1.5}%
\begin{tabular}{|>{\centering\arraybackslash}m{2.8cm}|>{\centering\arraybackslash}m{5.6cm}|>{\centering\arraybackslash}m{4.4cm}|}
\hline
\textbf{Naam} & \textbf{Kenmerk} & \textbf{Voorbeeld} \\ \hline
\textbf{\textcolor{accent}{vierkante matrix}} & evenveel rijen als kolommen: dimensie $\formule{n \times n}$ & $\begin{bmatrix} 2 & -1 \\ 0 & 5 \end{bmatrix}$ \\[4mm] \hline
\textbf{\textcolor{accent}{rijmatrix}} & één rij: dimensie $\formule{1 \times n}$ & $\begin{bmatrix} 4 & 0 & -3 \end{bmatrix}$ \\[2mm] \hline
\textbf{\textcolor{accent}{kolommatrix}} & één kolom: dimensie $\formule{m \times 1}$ & $\begin{bmatrix} 7 \\ 2 \end{bmatrix}$ \\[4mm] \hline
\textbf{\textcolor{accent}{nulmatrix}} $O$ & alle elementen zijn 0 & $\begin{bmatrix} 0 & 0 & 0 \\ 0 & 0 & 0 \end{bmatrix}$ \\[4mm] \hline
\end{tabular}\par}
\end{theorie}

Een nulmatrix bestaat in elke dimensie: welke je bedoelt, volgt uit de berekening. Een rijmatrix en een kolommatrix met dezelfde getallen zijn niet dezelfde matrix, want hun dimensie verschilt.

\subsection{Gelijke matrices}

Twee tabellen zijn pas dezelfde als ze even groot zijn en op elke plaats hetzelfde getal hebben. Voor matrices spreken we dat zo af.

\begin{theorie}
Twee matrices zijn \textbf{\textcolor{accent}{gelijk}} als ze dezelfde dimensie hebben en de elementen op dezelfde plaats gelijk zijn: \quad $\formule{A = B \iff a_{ij} = b_{ij}}$ voor elke $i$ en $j$.
\end{theorie}

\voorbeeld{}{Bepaal $x$ en $y$ zodat $\begin{bmatrix} 2x & 5 \\ y - 1 & 3 \end{bmatrix} = \begin{bmatrix} 8 & 5 \\ 4 & 3x - 9 \end{bmatrix}$.\label{vb:h1-gelijk}}{%
\uitleg{Een gelijkheid van matrices is in feite een reeks gelijkheden van getallen, één per plaats. We stellen ze op en lossen ze op. Elke vergelijking moet kloppen, ook die met een onbekende die we al kennen.}
\begin{stappen}
2x = 8 \ \Rightarrow \ x = 4 & De elementen in rij 1, kolom 1. \\[1ex]
y - 1 = 4 \ \Rightarrow \ y = 5 & De elementen in rij 2, kolom 1. \\[1ex]
3 = 3 \cdot 4 - 9 & Rij 2, kolom 2: ook die vergelijking moet kloppen. Met $x = 4$ staat er $3 = 3$, dus de oplossing is $x = 4$ en $y = 5$.
\end{stappen}}

Kwam er in de laatste stap iets anders uit dan $3 = 3$, dan waren er geen getallen $x$ en $y$ die de matrices gelijk maken.

\begin{oefeningen}
\oef{h1-gelijk}
\end{oefeningen}

\subsection{De getransponeerde matrix}

Je kunt de tabel van de filmzaal ook kantelen: de talen in de rijen en de dagen in de kolommen. De getallen blijven dezelfde, maar de eerste rij wordt de eerste kolom, de tweede rij de tweede kolom, enzovoort.

\begin{theorie}
De \textbf{\textcolor{accent}{getransponeerde}} $A\tr$ van een $m \times n$-matrix $A$ is de $n \times m$-matrix waarvan de rijen de kolommen van $A$ zijn:

\vspace{1mm}
\qquad $\formule{(A\tr)_{ij} = a_{ji}}$

\vspace{1mm}
Een vierkante matrix is \textbf{\textcolor{accent}{symmetrisch}} als $\formule{A\tr = A}$.
\end{theorie}

Je leest $A\tr$ als \textit{A getransponeerd}. Bij een symmetrische matrix is de hoofddiagonaal een spiegel: het element in rij $i$ en kolom $j$ is hetzelfde als dat in rij $j$ en kolom $i$.

\voorbeeld{}{Bepaal $A\tr$ voor $A = \begin{bmatrix} 1 & -2 & 0 \\ 4 & 3 & 6 \end{bmatrix}$.\label{vb:h1-transpose}}{%
\uitleg{De rijen van $A$ worden de kolommen van $A\tr$. Een $2 \times 3$-matrix wordt dus een $3 \times 2$-matrix.}
\begin{stappen}
A\tr = \begin{bmatrix} 1 & 4 \\ -2 & 3 \\ 0 & 6 \end{bmatrix} & De eerste rij $1,\ {-2},\ 0$ staat nu in de eerste kolom, de tweede rij $4,\ 3,\ 6$ in de tweede. Controle: $a_{13} = 0$ en $(A\tr)_{31} = 0$.
\end{stappen}}

\voorbeeld{}{Vier leerlingen werken het hele jaar in groepjes. De matrix $G$ geeft hoe vaak twee leerlingen samen in een groepje zaten. Leg uit waarom $G$ symmetrisch is.\label{vb:h1-symmetrisch}}{%
\uitleg{We kijken wat een element en het element in gespiegelde positie betekenen.}
\naast{\begin{minipage}{9.5cm}
Het element $g_{12}$ zegt hoe vaak leerling 1 met leerling 2 samenzat, en $g_{21}$ hoe vaak leerling 2 met leerling 1 samenzat. Dat is hetzelfde, dus $g_{12} = g_{21}$, en zo voor elk paar: $G\tr = G$. Op de hoofddiagonaal staan nullen, omdat een leerling niet met zichzelf in een groepje zit.
\end{minipage}}{$G = \begin{mat} 0 & 3 & 1 & 4 \\ 3 & 0 & 2 & 0 \\ 1 & 2 & 0 & 5 \\ 4 & 0 & 5 & 0 \end{mat}$}}

\nujij{Bepaal $B\tr$ voor de kolommatrix $B = \km{2 \\ 0 \\ -5}$.}{$B\tr = \begin{bmatrix} 2 & 0 & -5 \end{bmatrix}$: een rijmatrix.}

\begin{oefeningen}
\oef{h1-transpose}
\end{oefeningen}

%========================================================================================
\section{Optellen en vermenigvuldigen met een getal}

\subsection{De som van matrices}

In oefening \ref{oef:h1-verken} tel je de pagina's van september en oktober samen. Dat doe je plaats per plaats: de pagina's Frans voor het museum bij de pagina's Frans voor het museum. Dat lukt alleen als beide tabellen dezelfde vorm hebben.

\begin{theorie}
De \textbf{\textcolor{accent}{som}} van twee matrices met dezelfde dimensie krijg je door de elementen op dezelfde plaats op te tellen: \quad $\formule{(A + B)_{ij} = a_{ij} + b_{ij}}$

\vspace{1mm}
Het \textbf{\textcolor{accent}{verschil}} krijg je zo door af te trekken: \quad $\formule{(A - B)_{ij} = a_{ij} - b_{ij}}$
\end{theorie}

Matrices met een verschillende dimensie kun je niet optellen: $\km{1 & 2} + \km{1 & 2 & 3}$ bestaat niet.

\voorbeeld{}{Bereken $A + B$ en $A - B$ voor $A = \begin{bmatrix} 3 & -1 & 0 \\ 2 & 5 & -4 \end{bmatrix}$ en $B = \begin{bmatrix} -2 & 6 & 1 \\ 2 & -3 & 7 \end{bmatrix}$.\label{vb:h1-som}}{%
\uitleg{Beide matrices zijn $2 \times 3$, dus de som en het verschil bestaan, en ze zijn ook $2 \times 3$. We rekenen plaats per plaats, en letten op de mintekens.}
\begin{stappenbreed}
A + B = \begin{bmatrix} 3 + (-2) & -1 + 6 & 0 + 1 \\ 2 + 2 & 5 + (-3) & -4 + 7 \end{bmatrix} = \begin{bmatrix} 1 & 5 & 1 \\ 4 & 2 & 3 \end{bmatrix} & Elk element van $A$ plus het element van $B$ op dezelfde plaats. \\[4ex]
A - B = \begin{bmatrix} 3 - (-2) & -1 - 6 & 0 - 1 \\ 2 - 2 & 5 - (-3) & -4 - 7 \end{bmatrix} = \begin{bmatrix} 5 & -7 & -1 \\ 0 & 8 & -11 \end{bmatrix} & Nu trekken we af. Een negatief getal aftrekken is het tegengestelde optellen.
\end{stappenbreed}}

Optellen gebeurt plaats per plaats, en op elke plaats tel je gewone getallen op. Daarom erft de optelling van matrices de rekenregels die je al kent.

\begin{theorie}
\begin{stelling}{Eigenschappen van de optelling}
Voor matrices $A$, $B$ en $C$ met dezelfde dimensie:

\vspace{1mm}
{\centering\arrayrulecolor{accent}\renewcommand{\arraystretch}{1.6}%
\begin{tabular}{|>{\centering\arraybackslash}m{3.4cm}|>{\centering\arraybackslash$}m{6.4cm}<{$}|}
\hline
commutatief & \formule{A + B = B + A} \\ \hline
associatief & \formule{(A + B) + C = A + (B + C)} \\ \hline
nulmatrix & \formule{A + O = A} \\ \hline
tegengestelde & \formule{A + (-A) = O} \\ \hline
\end{tabular}\par}

\vspace{1mm}
Hierin is $-A$ de \textbf{\textcolor{accent}{tegengestelde matrix}}: elk element van teken veranderd.
\end{stelling}

\begin{bewijs}
We bewijzen de eerste regel; de andere gaan op dezelfde manier. Op elke plaats staat in $A + B$ en in $B + A$ een som van twee getallen, en voor getallen geldt
\bewijsregel{a_{ij} + b_{ij} = b_{ij} + a_{ij} .}
De matrices hebben dezelfde dimensie en op elke plaats hetzelfde element, dus ze zijn gelijk.
\end{bewijs}
\end{theorie}

\subsection{Een matrix vermenigvuldigen met een getal}

Stijgen alle prijzen van het vertaalbureau met 10~\%, dan vermenigvuldig je elke prijs met 1,1. Bij een matrix gaat dat net zo: elk element maal hetzelfde getal. Zo'n getal heet een \textbf{scalair}, om het te onderscheiden van een matrix.

\begin{theorie}
Het \textbf{\textcolor{accent}{product van een getal}} $r$ en een matrix $A$ krijg je door elk element met $r$ te vermenigvuldigen: \quad $\formule{(rA)_{ij} = r \cdot a_{ij}}$

\begin{stelling}{Rekenregels}
Voor getallen $r$, $s$ en matrices $A$, $B$ met dezelfde dimensie:

\vspace{1mm}
\qquad $\formule{r(A + B) = rA + rB}$ \qquad $\formule{(r + s)A = rA + sA}$ \qquad $\formule{r(sA) = (rs)A}$
\end{stelling}
\end{theorie}

Ook deze regels volgen uit de rekenregels voor getallen, plaats per plaats. Zo is $(-1)A = -A$, en \mbox{$A + A = 2A$}, zoals je zou verwachten.

\voorbeeld{}{Bereken $3A - 2B$ voor $A = \begin{bmatrix} 2 & 0 \\ -1 & 4 \end{bmatrix}$ en $B = \begin{bmatrix} 1 & -3 \\ 5 & 2 \end{bmatrix}$.\label{vb:h1-scalair}}{%
\uitleg{Zoals bij getallen gaat vermenigvuldigen voor aftrekken: eerst $3A$ en $2B$, daarna het verschil.}
\begin{stappenbreed}
3A = \begin{bmatrix} 6 & 0 \\ -3 & 12 \end{bmatrix}, \qquad 2B = \begin{bmatrix} 2 & -6 \\ 10 & 4 \end{bmatrix} & Elk element maal 3, en elk element maal 2. \\[4ex]
3A - 2B = \begin{bmatrix} 6 - 2 & 0 - (-6) \\ -3 - 10 & 12 - 4 \end{bmatrix} = \begin{bmatrix} 4 & 6 \\ -13 & 8 \end{bmatrix} & Plaats per plaats aftrekken.
\end{stappenbreed}}

Met die rekenregels kun je ook een vergelijking oplossen waarin de onbekende een matrix is. Je rekent zoals met getallen, alleen deel je niet door een matrix: je vermenigvuldigt met een getal.

\voorbeeld{}{Bepaal de matrix $X$ waarvoor $2X + A = B$, met $A = \begin{bmatrix} 1 & 5 \\ -2 & 0 \end{bmatrix}$ en $B = \begin{bmatrix} 7 & -1 \\ 4 & 6 \end{bmatrix}$.\label{vb:h1-vergelijking}}{%
\uitleg{We zonderen $X$ af zoals in een vergelijking met getallen: eerst $A$ naar de andere kant brengen, dan beide leden met $\tfrac12$ vermenigvuldigen. Pas op het einde vullen we de matrices in.}
\begin{stappen}
2X = B - A & Trek $A$ af van beide leden. \\[1ex]
X = \tfrac12 (B - A) & Vermenigvuldig beide leden met $\tfrac12$. \\[1ex]
X = \tfrac12 \begin{bmatrix} 6 & -6 \\ 6 & 6 \end{bmatrix} = \begin{bmatrix} 3 & -3 \\ 3 & 3 \end{bmatrix} & Nu pas rekenen: eerst $B - A$, dan de helft van elk element.
\end{stappen}
Controle: $2X + A = \km{6 & -6 \\ 6 & 6} + \km{1 & 5 \\ -2 & 0} = \km{7 & -1 \\ 4 & 6} = B$.}

\begin{werkwijze}
Breng alle termen met $X$ naar het ene lid en de rest naar het andere, zoals bij een vergelijking met getallen. Vermenigvuldig dan met het omgekeerde van het getal voor $X$. Vul de matrices pas op het einde in, en controleer door $X$ in te vullen.
\end{werkwijze}

\nujij{Bepaal $X$ als $X - \km{2 & 1 \\ 0 & 3} = 3\km{1 & 0 \\ 1 & 1}$.}{$X = \km{3 & 0 \\ 3 & 3} + \km{2 & 1 \\ 0 & 3} = \km{5 & 1 \\ 3 & 6}$.}

\begin{oefeningen}
\oef{h1-som}

\oef{h1-vergelijking}
\end{oefeningen}

%========================================================================================
\section{Het product van matrices}

\subsection{Een rij maal een kolom}

In oefening \ref{oef:h1-verken} bereken je wat het museum betaalt: elk aantal pagina's maal de prijs per pagina in die taal, en alles opgeteld. De aantallen staan in een rij, de prijzen in een kolom. Die berekening noemen we het product van de rij en de kolom.

\begin{theorie}
Het \textbf{\textcolor{accent}{product van een rij en een kolom}} met evenveel elementen is de $1 \times 1$-matrix met de som van de producten:

\vspace{1mm}
\qquad $\formule{\begin{bmatrix} a_1 & a_2 & \cdots & a_n \end{bmatrix} \cdot \begin{bmatrix} b_1 \\ b_2 \\ \vdots \\ b_n \end{bmatrix} = \begin{bmatrix} a_1 b_1 + a_2 b_2 + \cdots + a_n b_n \end{bmatrix}}$
\end{theorie}

\voorbeeld{}{Het museum liet in september 12 pagina's naar het Frans vertalen, 20 naar het Engels en 5 naar het Duits. Een pagina kost 25 euro voor het Frans, 20 euro voor het Engels en 30 euro voor het Duits. Schrijf de factuur als product van een rij en een kolom.\label{vb:h1-rijkolom}}{%
\uitleg{De aantallen komen in een rij, de prijzen in een kolom, in dezelfde volgorde van de talen. Dan staat elk aantal tegenover de juiste prijs.}
\begin{stappenbreed}
\begin{bmatrix} 12 & 20 & 5 \end{bmatrix} \cdot \begin{bmatrix} 25 \\ 20 \\ 30 \end{bmatrix} = \begin{bmatrix} 12 \cdot 25 + 20 \cdot 20 + 5 \cdot 30 \end{bmatrix} & Het eerste element van de rij maal het eerste van de kolom, het tweede maal het tweede, enzovoort. \\[4ex]
= \begin{bmatrix} 300 + 400 + 150 \end{bmatrix} = \begin{bmatrix} 850 \end{bmatrix} & Het museum betaalt 850 euro.
\end{stappenbreed}}

\subsection{Het product van twee matrices}

Het vertaalbureau heeft meer klanten, en ook een duurder tarief voor teksten die met spoed klaar moeten. We zetten de klanten onder elkaar en de tarieven naast elkaar. Elke klant met elk tarief is dan telkens een rij maal een kolom. Zo krijg je een hele tabel van producten.

\begin{theorie}
Is $A$ een $m \times n$-matrix en $B$ een $n \times p$-matrix, dan is hun \textbf{\textcolor{accent}{product}} $AB$ de $m \times p$-matrix waarvan het element in rij $i$ en kolom $j$ het product is van rij $i$ van $A$ en kolom $j$ van $B$:

\vspace{1mm}
\qquad $\formule{(AB)_{ij} = a_{i1} b_{1j} + a_{i2} b_{2j} + \cdots + a_{in} b_{nj}}$

\vspace{2mm}
Het product bestaat alleen als $A$ evenveel \emph{kolommen} heeft als $B$ \emph{rijen}: \quad $\formule{(m \times n) \cdot (n \times p) = m \times p}$
\end{theorie}

\figuren{\begin{tikzpicture}[font=\small]
  \matrix[getallen] (a) {2 \& -1 \& 3 \\ 0 \& 4 \& 1 \\};
  \matrix[getallen,right=9mm of a] (b) {1 \& 2 \\ 3 \& 0 \\ -2 \& 5 \\};
  \matrix[getallen,right=9mm of b] (c) {-7 \& 19 \\ 10 \& 5 \\};
  \node at ($(a.east)!.5!(b.west)$) {$\cdot$};
  \node at ($(b.east)!.5!(c.west)$) {$=$};
  \markeer{a-1-1}{a-1-3}
  \markeer{b-1-2}{b-3-2}
  \markeer{c-1-2}{c-1-2}
  \node[below=1mm of a,font=\scriptsize] {$2 \times 3$};
  \node[below=1mm of b,font=\scriptsize] {$3 \times 2$};
  \node[below=1mm of c,font=\scriptsize] {$2 \times 2$};
  \node[right=4mm of c,font=\scriptsize,align=left] {rij 1 maal kolom 2:\\$2 \cdot 2 + (-1) \cdot 0 + 3 \cdot 5 = 19$};
\end{tikzpicture}}

De figuur toont hoe één element ontstaat: de rij van $A$ schuift als het ware over de kolom van $B$. In het schema $(m \times n) \cdot (n \times p)$ moeten de twee binnenste getallen gelijk zijn, en de buitenste geven de dimensie van het product.

\voorbeeld{}{Bereken $AB$ voor $A = \begin{bmatrix} 2 & -1 & 3 \\ 0 & 4 & 1 \end{bmatrix}$ en $B = \begin{bmatrix} 1 & 2 \\ 3 & 0 \\ -2 & 5 \end{bmatrix}$.\label{vb:h1-product}}{%
\uitleg{Eerst kijken we of het product bestaat: $A$ heeft 3 kolommen en $B$ heeft 3 rijen, dus dat lukt, en het product is $2 \times 2$. Dan berekenen we de vier elementen, telkens een rij van $A$ maal een kolom van $B$.}
\begin{stappen}
(AB)_{11} = 2 \cdot 1 + (-1) \cdot 3 + 3 \cdot (-2) = -7 & Rij 1 van $A$ maal kolom 1 van $B$. \\[1ex]
(AB)_{12} = 2 \cdot 2 + (-1) \cdot 0 + 3 \cdot 5 = 19 & Rij 1 maal kolom 2: het element uit de figuur. \\[1ex]
(AB)_{21} = 0 \cdot 1 + 4 \cdot 3 + 1 \cdot (-2) = 10 & Rij 2 maal kolom 1. \\[1ex]
(AB)_{22} = 0 \cdot 2 + 4 \cdot 0 + 1 \cdot 5 = 5 & Rij 2 maal kolom 2. \\[1.5ex]
AB = \begin{bmatrix} -7 & 19 \\ 10 & 5 \end{bmatrix} & Elk element komt op de plaats van zijn rij en kolom.
\end{stappen}}

\begin{werkwijze}
Controleer met de dimensies of het product bestaat, en bepaal de dimensie van het resultaat. Voor het element in rij $i$ en kolom $j$: neem rij $i$ van de eerste matrix en kolom $j$ van de tweede, vermenigvuldig ze element per element en tel op.
\end{werkwijze}

\nujij{Bereken $\km{1 & 3 \\ -2 & 0} \cdot \km{4 \\ -1}$.}{Een $2 \times 2$ maal een $2 \times 1$ geeft een $2 \times 1$: $\km{4 - 3 \\ -8 + 0} = \km{1 \\ -8}$.}

\begin{oefeningen}
\oef{h1-product}
\end{oefeningen}

Bij een product in een context vertellen de opschriften wat je mag vermenigvuldigen. Een rij van de eerste matrix en een kolom van de tweede moeten over hetzelfde gaan: hier de talen.

\voorbeeld{}{De matrix $S$ geeft het aantal pagina's dat twee klanten in september lieten vertalen. De matrix $P$ geeft de prijs per pagina in euro, voor het gewone tarief en voor spoed. Bereken $SP$ en leg uit wat het element in rij 2, kolom 2 betekent.
\figuren{\small $S$:\ \begin{tabular}[t]{l|ccc} & Frans & Engels & Duits \\ \hline museum & 12 & 20 & 5 \\ reisbureau & 30 & 16 & 10 \end{tabular}
\qquad\qquad
$P$:\ \begin{tabular}[t]{l|cc} & gewoon & spoed \\ \hline Frans & 25 & 35 \\ Engels & 20 & 28 \\ Duits & 30 & 42 \end{tabular}}\label{vb:h1-context}}{%
\uitleg{De kolommen van $S$ zijn de talen, en de rijen van $P$ ook: in die volgorde mogen we vermenigvuldigen. Het product heeft de rijen van $S$ (de klanten) en de kolommen van $P$ (de tarieven).}
\begin{stappenbreed}
SP = \begin{bmatrix} 12 \cdot 25 + 20 \cdot 20 + 5 \cdot 30 & 12 \cdot 35 + 20 \cdot 28 + 5 \cdot 42 \\ 30 \cdot 25 + 16 \cdot 20 + 10 \cdot 30 & 30 \cdot 35 + 16 \cdot 28 + 10 \cdot 42 \end{bmatrix} & Elke klant maal elk tarief. \\[4ex]
SP = \begin{bmatrix} 850 & 1190 \\ 1370 & 1918 \end{bmatrix} & Een $2 \times 3$ maal een $3 \times 2$ geeft een $2 \times 2$-matrix. \\[3ex]
(SP)_{22} = 1918 & Rij 2 is het reisbureau, kolom 2 het spoedtarief: als alles met spoed moest, betaalde het reisbureau 1918 euro.
\end{stappenbreed}}

Het product $PS$ bestaat ook: een $3 \times 2$ maal een $2 \times 3$. Maar dan vermenigvuldig je een tarief met een klant, en dat betekent niets. Het product bestaan is dus niet genoeg: in een context moet het ook zinvol zijn.

\begin{werkwijze}
Zet de matrices zo dat de opschriften van de kolommen van de eerste matrix dezelfde zijn als die van de rijen van de tweede, in dezelfde volgorde. Het product krijgt dan de rijen van de eerste en de kolommen van de tweede.
\end{werkwijze}

\begin{oefeningen}
\oef{h1-context}
\end{oefeningen}

\subsection{Eigenschappen van het product}

Bij getallen maakt de volgorde van een product niet uit: $3 \cdot 5 = 5 \cdot 3$. Bij matrices hebben we net gezien dat $PS$ en $SP$ zelfs een andere dimensie kunnen hebben. Maar wat als beide producten dezelfde dimensie hebben? Probeer het eerst zelf.

\begin{oefeningen}
\oef{h1-commutatief}
\end{oefeningen}

Je vindt twee verschillende matrices: $AB \neq BA$. Het product van matrices is dus \emph{niet commutatief}. Bij een product van matrices moet je altijd zeggen in welke volgorde je vermenigvuldigt: $AB$ is \textit{A maal B}, met $A$ links. Andere regels voor getallen blijven wel gelden. Dat kun je nagaan op voorbeelden, en we nemen ze aan zonder bewijs.

\begin{theorie}
\begin{stelling}{Eigenschappen van het product}
Voor matrices waarvan de producten bestaan, en een getal $r$:

\vspace{1mm}
{\centering\arrayrulecolor{accent}\renewcommand{\arraystretch}{1.6}%
\begin{tabular}{|>{\centering\arraybackslash}m{3.4cm}|>{\centering\arraybackslash$}m{8.6cm}<{$}|}
\hline
associatief & \formule{(AB)C = A(BC)} \\ \hline
distributief & \formule{A(B + C) = AB + AC} \quad \text{en} \quad \formule{(A + B)C = AC + BC} \\ \hline
met een getal & \formule{r(AB) = (rA)B = A(rB)} \\ \hline
\end{tabular}\par}

\vspace{1mm}
In het algemeen is \ $\formule{AB \neq BA}$.
\end{stelling}
\end{theorie}

Omdat het product associatief is, mag je de haakjes weglaten: $ABC$ betekent $(AB)C$ en dat is ook $A(BC)$. De volgorde van de matrices mag je niet veranderen.

Bij getallen speelt 1 een bijzondere rol: $1 \cdot a = a \cdot 1 = a$. Bij matrices doet de eenheidsmatrix dat.

\begin{theorie}
De \textbf{\textcolor{accent}{eenheidsmatrix}} $I_n$ is de vierkante $n \times n$-matrix met 1 op de hoofddiagonaal en 0 op alle andere plaatsen:

\vspace{1mm}
\qquad $\formule{I_2 = \begin{bmatrix} 1 & 0 \\ 0 & 1 \end{bmatrix}} \qquad \formule{I_3 = \begin{bmatrix} 1 & 0 & 0 \\ 0 & 1 & 0 \\ 0 & 0 & 1 \end{bmatrix}}$

\begin{stelling}{Eenheidsmatrix}
Voor elke vierkante $n \times n$-matrix $A$ geldt \quad $\formule{A I_n = I_n A = A}$
\end{stelling}

\begin{bewijs}
We doen het voor $n = 2$; bij grotere matrices gaat het net zo. Neem $A = \km{a & b \\ c & d}$. Dan is
\bewijsregel{A I_2 = \begin{bmatrix} a \cdot 1 + b \cdot 0 & a \cdot 0 + b \cdot 1 \\ c \cdot 1 + d \cdot 0 & c \cdot 0 + d \cdot 1 \end{bmatrix} = \begin{bmatrix} a & b \\ c & d \end{bmatrix} = A ,}
en op dezelfde manier is $I_2 A = A$.
\end{bewijs}
\end{theorie}

Als het duidelijk is welke dimensie bedoeld is, schrijven we gewoon $I$. Een tweede verschil met getallen gaat over het getal 0. Als $a \cdot b = 0$, dan is $a = 0$ of $b = 0$. Bij matrices hoeft dat niet.

\voorbeeld{}{Bereken $AB$ voor $A = \begin{bmatrix} 1 & 2 \\ 2 & 4 \end{bmatrix}$ en $B = \begin{bmatrix} 2 & -4 \\ -1 & 2 \end{bmatrix}$. Wat valt je op?\label{vb:h1-nuldeler}}{%
\uitleg{We rekenen het product gewoon uit, en vergelijken met wat we van getallen gewoon zijn.}
\begin{stappen}
AB = \begin{bmatrix} 2 - 2 & -4 + 4 \\ 4 - 4 & -8 + 8 \end{bmatrix} = \begin{bmatrix} 0 & 0 \\ 0 & 0 \end{bmatrix} & Elk element is een rij van $A$ maal een kolom van $B$, en telkens heffen de twee termen elkaar op. \\[4ex]
AB = O, \text{ maar } A \neq O \text{ en } B \neq O & Het product van twee matrices kan de nulmatrix zijn, zonder dat een van beide de nulmatrix is.
\end{stappen}}

Daarom mag je in een gelijkheid van matrices niet zomaar \textit{schrappen}: uit $AB = AC$ volgt niet dat $B = C$.

\begin{oefeningen}
\oef{h1-eigenschap}
\end{oefeningen}

%========================================================================================
\section{Machten van een matrix}

Een getal tot de derde macht is drie keer dat getal maal zichzelf. Bij een matrix kan dat ook, maar alleen als ze vierkant is: een $2 \times 3$-matrix kun je niet met zichzelf vermenigvuldigen, want 3 kolommen passen niet op 2 rijen.

\begin{theorie}
Voor een vierkante matrix $A$ en een natuurlijk getal $n \geqslant 1$ is de \textbf{\textcolor{accent}{macht}}

\vspace{1mm}
\qquad $\formule{A^n = \underbrace{A \cdot A \cdots A}_{n \text{ factoren}}}$ \qquad en we spreken af dat $\formule{A^0 = I}$.
\end{theorie}

Omdat het product associatief is, maakt het niet uit welke factoren je eerst vermenigvuldigt. Zo is $A^3 = A^2 \cdot A = A \cdot A^2$, en $A^5 = A^2 \cdot A^3$. Zoals bij getallen geldt dus $A^m A^n = A^{m+n}$.

\voorbeeld{}{Bereken $A^2$ en $A^3$ voor $A = \begin{bmatrix} 1 & 2 \\ 0 & 1 \end{bmatrix}$, en zoek een formule voor $A^n$.\label{vb:h1-macht}}{%
\uitleg{We berekenen de eerste machten met de hand. Elke volgende macht is de vorige maal $A$. Dan kijken we welk element verandert, en hoe.}
\begin{stappen}
A^2 = \begin{bmatrix} 1 & 2 \\ 0 & 1 \end{bmatrix} \begin{bmatrix} 1 & 2 \\ 0 & 1 \end{bmatrix} = \begin{bmatrix} 1 & 4 \\ 0 & 1 \end{bmatrix} & Rij 1 maal kolom 2: $1 \cdot 2 + 2 \cdot 1 = 4$. De andere elementen blijven 1, 0 en 1. \\[4ex]
A^3 = A^2 \cdot A = \begin{bmatrix} 1 & 4 \\ 0 & 1 \end{bmatrix} \begin{bmatrix} 1 & 2 \\ 0 & 1 \end{bmatrix} = \begin{bmatrix} 1 & 6 \\ 0 & 1 \end{bmatrix} & Nu $1 \cdot 2 + 4 \cdot 1 = 6$: rechtsboven komt er telkens 2 bij. \\[4ex]
A^n = \begin{bmatrix} 1 & 2n \\ 0 & 1 \end{bmatrix} & Het element rechtsboven is 2, 4, 6, \ldots: twee keer de exponent.
\end{stappen}}

Let op: $A^2$ is niet de matrix met de kwadraten van de elementen. In het voorbeeld zou dat $\km{1 & 4 \\ 0 & 1}$ geven, toevallig hetzelfde, maar bij $\km{1 & 2 \\ 3 & 4}$ is $A^2 = \km{7 & 10 \\ 15 & 22}$ en niet $\km{1 & 4 \\ 9 & 16}$.

Ook de merkwaardige producten veranderen, omdat $AB$ en $BA$ verschillen. Met de distributiviteit krijg je
\[ (A + B)^2 = (A + B)(A + B) = A^2 + AB + BA + B^2 . \]
Alleen als $AB = BA$ mag je de twee middelste termen samennemen tot $2AB$.

Hogere machten met de hand uitrekenen is veel werk. Je rekentoestel kan het voor je doen.

\voorbeeld{}{Bereken $A^{10}$ voor $A = \begin{bmatrix} 0{,}5 & 1 \\ 0{,}5 & 0 \end{bmatrix}$ met je rekentoestel.\label{vb:h1-rekentoestel}}{%
\uitleg{Tien keer met de hand vermenigvuldigen is te veel werk. We voeren de matrix in op het rekentoestel en laten het de macht berekenen.}
\begin{stappen}
A = \begin{bmatrix} 0{,}5 & 1 \\ 0{,}5 & 0 \end{bmatrix} & Voer de matrix in via het matrixmenu: eerst de dimensie $2 \times 2$, dan de elementen rij per rij. \\[4ex]
A^{10} \approx \begin{bmatrix} 0{,}667 & 0{,}666 \\ 0{,}333 & 0{,}334 \end{bmatrix} & Reken in het gewone scherm met de naam van de matrix, zoals met een getal: $A$ tot de tiende macht. Rond de elementen af op drie decimalen.
\end{stappen}}

De kolommen van $A^{10}$ zijn bijna gelijk. Bereken je $A^{20}$, dan zijn ze op drie decimalen gelijk. In hoofdstuk 2 zie je wat dat betekent.

\begin{werkwijze}
Een kleine macht van een kleine matrix reken je met de hand uit: telkens de vorige macht maal $A$. Zoek je een patroon, schrijf dan de eerste machten onder elkaar en kijk welk element hoe verandert. Voor een hoge macht gebruik je je rekentoestel.
\end{werkwijze}

\lkr{Op de TI-84: \texttt{2nd} \texttt{x}$^{-1}$ (MATRIX), tab EDIT om de matrix in te voeren, en in het rekenscherm \texttt{[A]\^{}10}. Leerlingen ronden soms af vóór ze verder rekenen: laat ze altijd met de matrix in het geheugen werken.}

\nujij{Bereken $A^2$ en $A^3$ voor $A = \km{2 & 0 \\ 0 & -1}$. Wat is $A^n$?}{$A^2 = \km{4 & 0 \\ 0 & 1}$, $A^3 = \km{8 & 0 \\ 0 & -1}$, en $A^n = \km{2^n & 0 \\ 0 & (-1)^n}$.}

\begin{oefeningen}
\oef{h1-macht}
\end{oefeningen}

%========================================================================================
\opdrachtenkop

\begin{oefeningen}
\oef{h1-o1}

\oef{h1-o2}

\oef{h1-o3}

\oef{h1-o4}

\oef{h1-o5}

\oef{h1-o6}

\oef{h1-o7}

\oef{h1-o8}

\oef{h1-o9}

\oef{h1-o10}

\needspace{18\baselineskip}
\oef{h1-o11}

\oef{h1-o12}

\oef{h1-o13}

\oef{h1-o14}

\oef{h1-o15}

\oef{h1-o16}

\oef{h1-o17}

\oef{h1-o18}

\oef{h1-o19}

\oef{h1-o20}

\oef{h1-o21}

\oef{h1-o22}

\oef{h1-o23}
\end{oefeningen}

%----------------------------------------------------------------------------------------
% KOPPEL: verbanden met andere leerstof
%----------------------------------------------------------------------------------------
\needspace{20\baselineskip}
\opdrachtsectie{Koppel}

\lkr{Spiegelingen en draaiingen kennen de leerlingen uit de 2de graad. Statistiek en Rijen staan in dezelfde graad: kijk na of ze al aan bod kwamen, want het pakket legt de volgorde van de cursussen niet vast.}

\begin{oefeningen}
\oef{h1-koppel-1}

\oef{h1-koppel-2}

\oef{h1-koppel-3}
\end{oefeningen}

%----------------------------------------------------------------------------------------
\opdrachtsectie{Test jezelf}

\begin{oefeningen}
\oef{h1-t1}

\oef{h1-t2}

\oef{h1-t3}

\oef{h1-t4}

\oef{h1-t5}

\oef{h1-t6}

\oef{h1-t7}

\oef{h1-t8}
\end{oefeningen}

\herhaalschema{%
  \fv{oef:h1-t1}{vb:h1-element}{oef:h1-element}%
  \fv{oef:h1-t2}{vb:h1-transpose}{oef:h1-transpose}%
  \fv{oef:h1-t3}{vb:h1-scalair}{oef:h1-som}%
  \fv{oef:h1-t4}{vb:h1-vergelijking}{oef:h1-vergelijking}%
  \fv{oef:h1-t5}{vb:h1-product}{oef:h1-product}%
  \fv{oef:h1-t6}{vb:h1-context}{oef:h1-context}%
  \fv{oef:h1-t7}{vb:h1-nuldeler}{oef:h1-eigenschap}%
  \fv{oef:h1-t8}{vb:h1-macht}{oef:h1-macht}}

\ifzelfstudie
\section*{Maak je eigen samenvatting}
Sluit dit boek en schrijf in je schrift: wat een matrix, een element en de dimensie zijn, de bijzondere matrices, wat de getransponeerde is, hoe je optelt en met een getal vermenigvuldigt, wanneer een product bestaat en hoe je het berekent, welke rekenregels voor getallen niet gelden bij matrices, en wat een macht is. Vergelijk daarna met de kaders.
\fi
